\section{Population, rules-dialect and arbitration results}
\label{app:dialects}


\subsection{The round-robin population matrix}
\label{app:pop-matrix}

Table~\ref{tab:elo} in \S\ref{sec:results-pop} summarises the round-robin as
Bradley--Terry ratings. Table~\ref{tab:matrix} gives the win rates the ratings
were fitted to, so that a reader can check any individual pairing rather than
taking the rating on trust --- which matters, because Bradley--Terry ratings are
not redundancy-invariant and cannot represent cycles \cite{balduzzi-reeval}.

\begin{table}[!htbp]
\centering
\caption{Primary head-to-head (E3). Population: \numHtHDeals{} held-out deals in
all six seat rotations, \numHtHGames{} games per row, seed $90210$, full rules
dialect, \vfast{} frozen at the shipping configuration before the bank was run,
against the eight-policy panel. Intervals: deal-clustered percentile bootstrap on
the per-deal win count, \numBootstrapDraws{} resamples. ``Mean half-suits'' is
half-suits won of nine; the two declaration columns count \emph{voluntary}
declarations only, forced endgame declarations being tallied separately in the
run artifact and not shown here (\S\ref{sec:protocol-metrics}). Action-limit rate
\numLimitRate{} on every row.}
\label{tab:h2h}
\small
\begin{tabular}{lrrrrrr}
\toprule
Opponent & Win rate & 95\% CI & Mean sets & Ask acc. & Decl. acc. & Decl./game \\
\midrule
FishBot v0.3 & 75.07\% & 73.7--76.4 & 5.457 & 55.81\% & 98.55\% & 4.80 \\
Turn-starvation lockout & 78.12\% & 76.9--79.4 & 5.542 & 48.01\% & 98.14\% & 5.16 \\
Posterior detective & 75.74\% & 74.5--77.0 & 5.440 & 52.79\% & 98.36\% & 5.12 \\
FishBot v0.2 & 84.24\% & 83.1--85.3 & 5.857 & 54.42\% & 98.85\% & 5.17 \\
Adaptive diversifier & 92.14\% & 91.3--93.0 & 6.442 & 64.05\% & 97.42\% & 6.14 \\
Focused hunter & 97.64\% & 97.2--98.1 & 7.000 & 55.69\% & 97.51\% & 5.09 \\
Misdirection artist & 99.95\% & 99.9--100.0 & 8.161 & 55.12\% & 98.16\% & 3.80 \\
Random legal control & 100.00\% & 100.0--100.0 & 8.579 & 54.26\% & 96.72\% & 6.99 \\
\bottomrule
\end{tabular}
\end{table}

\numWorstCase\% is the \emph{lowest} win rate \vfast{} recorded against any
member of this eight-policy panel, on this bank, under this rules dialect and
under the default arbitration order of \S\ref{sec:rules}, two alternatives to
which put it slightly lower (\S\ref{sec:robustness}). It is not a claim about
arbitrary opponents. Deals were held out from fitting and selection, opponent
identities mostly were not --- six of the eight panel members were also in the
fitting panel, only the misdirection artist and the random legal control were
withheld (\S\ref{sec:protocol-seeds}) --- so these experiments measure
generalisation to new deals and only weakly to new strategic populations. That
scoping governs every result in this section and is not repeated at each one. For
scale, \vthree{}'s worst published result against its three credible challengers
was \numVthreeWorstPublished\% \cite{fishbot03}.

The margin is concentrated in declaration reliability rather than in ask
conversion. Against the turn-starvation lockout \vfast{} converts the lowest
fraction of asks in the panel (\numVsLockoutAskAcc\%) and still wins by a wider
margin (\numVsLockout\%) than against the comparably strong \vthree{}
(\numVsVthree\%, ask accuracy \numVsVthreeAskAcc\%) and posterior detective
(\numVsDetective\%), while declaring correctly \numVsVthreeDeclAcc\% of the time
on games where \vthree{} manages \numVthreeDeclAccSame\%. Every incorrect
declaration awards the half-suit to the opposing team in this dialect, so a
policy that converts fewer asks can still win by giving away fewer half-suits.
This describes where the observed margin sits; it is not a causal decomposition,
and \S\ref{sec:ablations} reports what the frozen-policy ablations do and do not
establish about the components involved. \vfast{} also uses the out-of-turn rule
heavily, \numVsVthreeOOT{} declarations out of turn per game against \vthree{};
\S\ref{sec:robustness} reports how the margin moves without it.

\subsection{Population ratings}
\label{sec:results-pop}

\begin{table}[!htbp]
\centering
\caption{Bradley--Terry ratings (E4) in Elo with \vthree{} at zero. Population:
the round-robin of Table~\ref{tab:matrix} --- \numMatrixDeals{} deals in six seat
rotations, \numMatrixGamesPair{} games per ordered pair, seed $515151$, nine
policies, all weights frozen. The mean win-rate column averages over the row's
opponents. Ratings are point estimates fitted to the full matrix; no intervals
are constructed for them.}
\label{tab:elo}
\small
\begin{tabular}{lrr}
\toprule
Policy & Elo (v0.3 $=0$) & Mean win rate \\
\midrule
FishBot v0.4 & +216 & 87.73\% \\
Posterior detective & +8 & 68.84\% \\
Turn-starvation lockout & +5 & 68.48\% \\
FishBot v0.3 & +0 & 68.02\% \\
FishBot v0.2 & -32 & 64.74\% \\
Adaptive diversifier & -218 & 46.46\% \\
Focused hunter & -441 & 29.60\% \\
Random legal control & -759 & 13.26\% \\
Misdirection artist & -1029 & 2.86\% \\
\bottomrule
\end{tabular}
\end{table}

\begin{table}[!htbp]
\centering
\caption{Round-robin win rates (E4), row policy against column policy, per cent.
Population: \numMatrixDeals{} deals in all six seat rotations for every ordered
pair, seed $515151$, full rules dialect, all weights frozen. Each cell pools both
orderings, so a cell is \numMatrixGamesCell{} games, \numMatrixGamesPair{} per
ordered pair. Cells are point estimates; no intervals per cell.}
\label{tab:matrix}
\scriptsize
\setlength{\tabcolsep}{3pt}
\begin{tabular}{lrrrrrrrrr}
\toprule
Row vs column & \vfast{} & \vthree{} & \vtwo{} & Lockout & Detective & Diversifier & Hunter & Misdirection & Random \\
\midrule
\vfast{} & --- & 75.7 & 83.5 & 78.7 & 74.0 & 92.7 & 97.3 & 99.9 & 100.0 \\
\vthree{} & 24.3 & --- & 51.5 & 49.5 & 49.8 & 78.8 & 91.0 & 99.3 & 99.9 \\
\vtwo{} & 16.5 & 48.5 & --- & 41.6 & 44.8 & 76.0 & 91.0 & 99.6 & 99.9 \\
Turn-starvation lockout & 21.3 & 50.5 & 58.4 & --- & 49.9 & 77.0 & 91.2 & 99.5 & 100.0 \\
Posterior detective & 26.0 & 50.2 & 55.2 & 50.1 & --- & 76.9 & 92.8 & 99.5 & 100.0 \\
Adaptive diversifier & 7.3 & 21.2 & 24.0 & 23.0 & 23.1 & --- & 74.4 & 99.5 & 99.3 \\
Focused hunter & 2.7 & 9.0 & 9.0 & 8.8 & 7.2 & 25.6 & --- & 93.1 & 81.5 \\
Misdirection artist & 0.1 & 0.7 & 0.4 & 0.5 & 0.5 & 0.5 & 6.9 & --- & 13.4 \\
Random legal control & 0.0 & 0.1 & 0.1 & 0.0 & 0.0 & 0.7 & 18.5 & 86.6 & --- \\
\bottomrule
\end{tabular}
\end{table}


\subsection{Validating the ported \vthree{} baseline}
\label{app:port}

Every comparison in this paper is against a re-implementation of \vthree{} in the
v0.4 engine, so that re-implementation has to be shown faithful before its scores
mean anything. Table~\ref{tab:port} replays it under the v0.3 rules dialect on
the v0.3 seed bank and compares against the published figures.

\begin{table}[!htbp]
\centering
\caption{\vthree{} as published \cite{fishbot03} versus the C++ port. Population:
\numPortDeals{} deals in two orientations, \numPortGames{} games per row, seed
$20260820$, v0.3 dialect, the seven opponents of the v0.3 study, ported weights
frozen. Intervals: deal-clustered percentile bootstrap on the port's win rate,
\numBootstrapDraws{} resamples.}
\label{tab:port}
\small
\begin{tabular}{lrrr}
\toprule
Opponent & Published v0.3 & C++ port & 95\% CI \\
\midrule
Turn-starvation lockout & 57.85\% & 57.60\% & 55.5--59.8 \\
Posterior detective & 57.20\% & 59.10\% & 57.0--61.3 \\
FishBot v0.2 & 56.50\% & 56.60\% & 54.5--58.7 \\
Adaptive diversifier & 86.15\% & 84.75\% & 83.2--86.3 \\
Focused hunter & 95.15\% & 94.55\% & 93.5--95.5 \\
Misdirection artist & 99.20\% & 98.70\% & 98.2--99.2 \\
Random legal control & 100.00\% & 100.00\% & 100.0--100.0 \\
\bottomrule
\end{tabular}
\end{table}

The engine implements every consequential rule choice of \S\ref{sec:rules} as a
switch, so the same frozen \vfast{} configuration can be replayed under other
dialects and under other resolutions of simultaneous voluntary declarations
without recompiling or refitting. This appendix reports both sweeps. Neither is
a claim about which dialect is correct; both are sensitivity analyses on the
modelling choices the main results are conditioned on.

\subsection{Rules dialects}
\label{app:dialects-rules}

\begin{table}[ht]
\centering
\caption{E7: the frozen \vfast{} configuration (\texttt{v04:mgate=0.008}, the
\numParams{}-coordinate fitted vector unchanged, no refitting) replayed under
three rules dialects against three opponents. Each cell is \numEsevenDeals{}
deals played in six seat rotations, \numEsevenGames{} games; seeds 828282 (v0.3
dialect), 838383 (48-card, eight-half-suit form) and 848484 (out-of-turn
declarations disabled), all disjoint from the fitting and selection banks.
``Win rate'' is \vfast{}'s share of games won; intervals are deal-clustered
percentile bootstrap, the 2.5\% and 97.5\% percentiles of \numBootstrapDraws{}
resamples of deals. Every row recorded an action-limit rate of \numLimitRate.
The full-dialect comparators quoted in the text come from E3
(Table~\ref{tab:h2h}), which is a different bank (\numHtHDeals{} deals, seed
90210), so the comparison is across banks as well as across dialects.}
\label{tab:rules}
\small
\begin{tabular}{llrr}
\toprule
Rule dialect & Opponent & Win rate & 95\% CI \\
\midrule
v0.3 dialect (turn-only declarations) & FishBot v0.3 & 64.46\% & 62.6--66.4 \\
48-card, eight half-suits & FishBot v0.3 & 72.54\% & 70.8--74.2 \\
No out-of-turn declarations & FishBot v0.3 & 67.71\% & 65.9--69.5 \\
v0.3 dialect (turn-only declarations) & Turn-starvation lockout & 77.25\% & 75.6--78.9 \\
48-card, eight half-suits & Turn-starvation lockout & 74.79\% & 73.2--76.4 \\
No out-of-turn declarations & Turn-starvation lockout & 80.17\% & 78.5--81.8 \\
v0.3 dialect (turn-only declarations) & Posterior detective & 76.21\% & 74.4--78.0 \\
48-card, eight half-suits & Posterior detective & 72.21\% & 70.5--73.9 \\
No out-of-turn declarations & Posterior detective & 74.75\% & 73.0--76.5 \\
\bottomrule
\end{tabular}
\end{table}

Table~\ref{tab:rules} shows that the size of the margin depends on the dialect,
and that the dependence is not in the same direction for every opponent.

Against \vthree{}, the margin is most sensitive to the declaration timing rules.
Under the v0.3 dialect, in which a player may declare only on their own turn,
\vfast{} wins \numEsevenVthreeLegacy\% (95\% deal-clustered percentile-bootstrap
interval \numEsevenVthreeLegacyCI). Under the full dialect with out-of-turn
declarations disabled but everything else intact it wins
\numEsevenVthreeNoOOT\% (\numEsevenVthreeNoOOTCI). Both are below the
\numVsVthree\% recorded under the full dialect in E3. The 48-card,
eight-half-suit form gives \numEsevenVthreeEight\% (\numEsevenVthreeEightCI).

Against the turn-starvation lockout the ordering is different. Disabling
out-of-turn declarations raises the margin to \numEsevenLockoutNoOOT\%, above
the \numVsLockout\% of E3, while the v0.3 dialect gives
\numEsevenLockoutLegacy\% and the 48-card form gives \numEsevenLockoutEight\%,
the lowest of the three. Against the posterior detective the v0.3 dialect gives
the highest of the three figures, \numEsevenDetectiveLegacy\%, against
\numVsDetective\% under the full dialect in E3; disabling out-of-turn
declarations gives \numEsevenDetectiveNoOOT\% and the 48-card form
\numEsevenDetectiveEight\%. The interval for each of these six cells is in
Table~\ref{tab:rules} and is a deal-clustered percentile bootstrap on the same
construction.

Two quantities do not move across the sweep. No cell in E7 recorded an
action-limit game: on these three banks, under the graduated forcing rule of
\S\ref{sec:locked}, every game terminated through play. That is a measurement on
\numEsevenGames{} games per cell, not a termination guarantee, and it does not
extend to dialects or configurations outside the three tested here
(\S\ref{sec:termination}). And \vfast{} retains a substantial margin in every
cell: the smallest figure in the table is \numEsevenVthreeLegacy\%. The dialect
changes the size of the margin, not its sign, over the three dialects and three
opponents tested. We do not attempt to explain the opponent-specific orderings;
disentangling them would require an experiment that varies one rule at a time
against a policy whose response to that rule is independently characterised,
which E7 does not do.

\subsection{Simultaneous-declaration arbitration}
\label{app:dialects-arb}

When more than one player would offer a declaration at the same declaration
opportunity, the driver must choose one. The shipped choice is the lowest seat
index. This is a modelling decision rather than a rule of the game, and the
main results are conditioned on it, so E17 replays the primary bank under all
three orders the engine supports: lowest seat index (\texttt{-{}-arb=low}, the
default), highest seat index (\texttt{-{}-arb=high}) and a scan beginning from
the current turn-holder (\texttt{-{}-arb=turn}).

\begin{table}[ht]
\centering
\caption{E17: sensitivity of the primary head-to-head results to the order in
which simultaneous voluntary declarations are arbitrated. The frozen \vfast{}
configuration (\texttt{v04:mgate=0.008}, no refitting) replayed against three
opponents under three arbitration orders; \numEseventeenDeals{} deals played in
six seat rotations, \numHtHGames{} games per cell, seed 90210, full dialect.
``Win rate'' is \vfast{}'s share of games won; intervals are deal-clustered
percentile bootstrap, the 2.5\% and 97.5\% percentiles of \numBootstrapDraws{}
resamples of deals. The last column is \vfast{}'s out-of-turn declarations per
game. The lowest-seat rows are the shipped configuration and coincide with
Table~\ref{tab:h2h}.}
\label{tab:arbitration}
\small
\begin{tabular}{llrrr}
\toprule
Arbitration order & Opponent & Win rate & 95\% CI & Out-of-turn decl./game \\
\midrule
Lowest seat index (default) & \vthree{} & 75.07\% & 73.71--76.40 & 3.43 \\
 & Turn-starvation lockout & 78.12\% & 76.88--79.36 & 3.59 \\
 & Posterior detective & 75.74\% & 74.45--77.02 & 3.57 \\
Highest seat index & \vthree{} & 74.93\% & 73.60--76.26 & 3.46 \\
 & Turn-starvation lockout & 78.05\% & 76.81--79.26 & 3.58 \\
 & Posterior detective & 75.74\% & 74.43--77.02 & 3.59 \\
Scan from the current turn-holder & \vthree{} & 74.38\% & 73.02--75.74 & 2.06 \\
 & Turn-starvation lockout & 78.19\% & 76.95--79.43 & 2.12 \\
 & Posterior detective & 75.76\% & 74.48--77.05 & 2.11 \\
\bottomrule
\end{tabular}
\end{table}

Across the nine cells of Table~\ref{tab:arbitration} the win rate moves by at
most \numArbMaxMove{} percentage points, and every deal-clustered
percentile-bootstrap interval overlaps the corresponding intervals under the
other two orders. Against \vthree{} the three orders give \numVsVthree\%,
\numArbHighVthree\% and \numArbTurnVthree\%; against the lockout,
\numVsLockout\%, \numArbHighLockout\% and \numArbTurnLockout\%; against the
detective, \numVsDetective\%, \numArbHighDetective\% and
\numArbTurnDetective\%. Declarations per game are unchanged to the reported
precision under all three orders (Table~\ref{tab:arbitration}).

The one quantity that does move is the rate of out-of-turn declarations. Under
either seat-index order it is between \numArbOOTSeatLo{} and
\numArbOOTSeatHi{} per game; under the scan from the turn-holder it falls to
between \numArbOOTTurnLo{} and \numArbOOTTurnHi{}. That is the expected
consequence of the change --- the turn-holder is examined first and, when it
also offers a declaration, resolves the opportunity in turn rather than out of
turn --- and it is a change in how declarations are labelled at least as much as
a change in which declarations happen, since the total declaration rate is
unchanged.

The reported margins are therefore not an artifact of the arbitration order.
Two limits on that statement should be kept in view. Only three orders were
tested, all of them deterministic functions of seat index or turn position; a
randomised tie-break, or an order that depends on the game state in some other
way, was not run. And one candidate order was deliberately excluded: ranking
simultaneous declarers by their stated confidence in their own allocation. That
rule would compare private quantities across seats, so it would let the
arbitration mechanism act as a channel through which one player's belief
influences the outcome of another player's decision, and the resulting simulator
would no longer satisfy the information-safety property audited in E1.
